\clearpage
\section{Lösungen}
\label{bin:solutions}
Ein Zählbeweis erklärt, welche Wege zusammengefasst werden und warum
jeder Weg genau einmal vorkommt. Die folgenden Lösungen führen auch
die Umkehrung der Zuordnungen und die nötigen Randfälle aus.

\Needspace{7\baselineskip}
\subsection*{Zu Aufgabe 1: Wege und Eltern}
\textbf{(a)} Nach $(4,2)$ führen genau die sechs Folgen
\[
 RRLL,\quad RLRL,\quad RLLR,\quad LRRL,\quad LRLR,\quad LLRR.
\]
Ihre $R$-Positionen sind der Reihe nach
\[
 \{1,2\},\ \{1,3\},\ \{1,4\},\ \{2,3\},\ \{2,4\},\ \{3,4\}.
\]
Dies sind alle Paare aus vier Positionen. Jedes Paar legt die beiden
$R$ und damit auch die beiden $L$ eindeutig fest; kein Weg fehlt oder
kommt doppelt vor. Zum Einzeichnen gehst du zum Beispiel für $RLLR$
von $(0,0)$ über $(1,1),(2,1),(3,1)$ nach $(4,2)$.

\textbf{(b)} Auf $R$ enden $RLLR,LRLR,LLRR$: Vorher liegen drei Schritte
mit einem $R$ vor, also die Stelle $(3,1)$. Auf $L$ enden
$RRLL,RLRL,LRRL$: Sie kommen von $(3,2)$. Jede Klasse hat drei Wege;
daher $\binom42=\binom31+\binom32=3+3=6$.

\textbf{(c)} Vor dem letzten $R$ wurden erst $k-1$ der insgesamt $k$
Rechtsschritte ausgeführt; vor einem letzten $L$ schon alle $k$.
Der obere Index ist vorher jeweils $n$. Entfernen und Wiederanhängen
des letzten Schritts sind umkehrbar. Die beiden Klassen sind getrennt
und vollständig, also
\[
 \binom{n+1}k=\binom n{k-1}+\binom nk.
\]
Bei $k=n+1$ ist der zweite Beitrag $0$: In $n$ Schritten sind $n+1$
Rechtsschritte unmöglich. Bei $k=0$ gibt es nur den Weg aus lauter $L$;
seine Anzahl ist $1$, ein Index $-1$ wird nicht benötigt.

\textbf{(d)} Gezählt werden Wege mit jeweils genau zwei $R$ und
insgesamt $2$, $3$ beziehungsweise $4$ Schritten: Es sind $1$, $3$
und $6$ Wege. Bei festem $k=2$ ist jede nächste Stelle über einen
Schritt $L$ erreichbar; deshalb verläuft diese Diagonale schräg nach
links unten. Der Stellenindex gehört zur jeweiligen Zeile.

\Needspace{7\baselineskip}
\subsection*{Zu Aufgabe 2: Ränder und Spiegelung}
\textbf{(a)} Die Werte sind in der angegebenen Reihenfolge
\[
 1,\quad0,\quad1,\quad9,\quad9,\quad1,\quad0.
\]
Für $\binom00$ zählt der leere Weg. Sieben $R$ bei null Schritten
und zehn $R$ bei neun Schritten sind unmöglich. Nach $(9,0)$ führt
nur $LLLLLLLLL$, nach $(9,9)$ nur $RRRRRRRRR$. Bei $(9,1)$ kann
das einzige $R$ an jeder der neun Positionen stehen; bei $(9,8)$
gilt dasselbe für das einzige $L$.

\textbf{(b)} Aus $k$ Rechtsschritten und $n-k$ Linksschritten werden
durch Vertauschen $n-k$ Rechtsschritte und $k$ Linksschritte.
Der Spiegelweg endet bei $(n,n-k)$. Erneutes Vertauschen stellt den
ursprünglichen Weg wieder her. Somit ist die Zuordnung in beide
Richtungen eindeutig, und $\binom nk=\binom n{n-k}$.
Auch der leere Weg wird auf sich selbst abgebildet.

\textbf{(c)} Die Symmetrie setzt $0\le k\le n$ voraus. Hier ist
$5>3$, und $3-5$ ist keine zugelassene natürliche Differenz.
Fünf Rechtsschritte in nur drei Schritten sind unmöglich; also $\binom35=0$.

\Needspace{7\baselineskip}
\subsection*{Zu Aufgabe 3: Rechtsschrittpositionen}
\textbf{(a)} Wähle drei der zwölf Positionen für $R$; die übrigen sind
$L$. Umgekehrt liefert jeder Weg genau seine drei $R$-Positionen.
Die Anzahl der Wege und der Auswahlen ist deshalb dieselbe. Es gilt
\[
 3!\,9!\binom{12}3=12!=12\cdot11\cdot10\cdot9!.
\]
Weil $9!>0$ ist, dürfen wir kürzen. Mit $3!=6$ ergibt sich
$6\binom{12}3=1320$ und damit $\binom{12}3=220$.

\textbf{(b)} Die Spiegelung vertauscht drei $R$ mit neun $R$.
Daher gibt es nach $(12,9)$ ebenfalls $\binom{12}9=\binom{12}3=220$ Wege.

\textbf{(c)} Wählen wir zuerst eine Position und danach eine andere,
haben wir $n(n-1)$ Möglichkeiten. Ein Weg mit zwei $R$ entsteht
zweimal: Seine beiden $R$-Positionen können in beiden Reihenfolgen
gewählt werden. Weitere Möglichkeiten für denselben Weg gibt es nicht.
Also gilt $2\binom n2=n(n-1)$.

Die Fakultätenformel bestätigt dies:
\[
 2!\,(n-2)!\binom n2=n!=n(n-1)(n-2)!.
\]
Der gemeinsame Faktor $(n-2)!$ ist positiv, auch bei $n=2$, wo er
$0!=1$ ist. Kürzen und $2!=2$ liefern dieselbe Gleichung.

\Needspace{7\baselineskip}
\subsection*{Zu Aufgabe 4: Ausmultiplizieren als Weg}
\textbf{(a)} Die Folge $RLLRL$ liefert $abbab=a^2b^3$. Genau zwei
der fünf Faktoren müssen $a$ liefern. Das entspricht den Wegen nach
$(5,2)$, also $\binom52=10$ Möglichkeiten. Ausmultiplizieren und
Zusammenfassen ergibt deshalb den Koeffizienten $10$.

\textbf{(b)} Die Wege mit $k$ Rechtsschritten liefern jeweils
$a^kb^{5-k}$. Zeile $5$ ergibt
\[
 \begin{aligned}
 (a+b)^5={}&b^5+5ab^4+10a^2b^3\\
           &+10a^3b^2+5a^4b+a^5.
 \end{aligned}
\]
Für $a=1$, $b=10$ folgt
\[
 \begin{aligned}
 11^5&=100000+50000+10000+1000+50+1\\
     &=161051.
 \end{aligned}
\]

\textbf{(c)} Bei $n=0$ gibt es keine Faktoren zur Auswahl und genau
einen leeren Weg. Sein leeres Produkt hat den Wert $1$. Links ist
$(a+b)^0=1$, rechts steht nur $\binom00a^0b^0=1$.
Auch für $a=b=0$ gilt somit $(0+0)^0=1=\binom00\,0^0\,0^0$,
mit der vereinbarten Konvention $0^0=1$.

\Needspace{7\baselineskip}
\subsection*{Zu Aufgabe 5: Endpunkte und Farben}
\textbf{(a)} Bei jedem der acht Schritte stehen unabhängig $L$ und
$R$ zur Wahl. Daher gibt es $2^8=256$ Wege. Jeder endet an genau einer
Stelle $(8,k)$; dorthin führen $\binom8k$ Wege. Die Endpunktklassen
sind getrennt und erfassen alle Wege. Also ist
\[
 \sum_{k=0}^8\binom8k=256.
\]

\textbf{(b)} Nur die beiden Wege aus lauter $L$ oder lauter $R$ enden
am Rand. Es bleiben $256-2=254$ Wege zu inneren Stellen.

\textbf{(c)} Bei jedem Schritt wählen wir jetzt zwischen
$L$, $R_{\mathrm{rot}}$ und $R_{\mathrm{blau}}$. Das gibt insgesamt
$3^n$ gefärbte Wege. Zu einem festen ungefärbten Weg nach $(n,k)$
gibt es dagegen $2^k$ Färbungen: Jede seiner $k$ Rechtsschrittpositionen
erhält unabhängig eine der beiden Farben.

Ein gefärbter Weg bestimmt eindeutig seinen ungefärbten Weg und die
Farben an dessen Rechtsschritten. Umgekehrt ergibt jede solche Kombination
einen gefärbten Weg. Daher enthält die Klasse mit Endpunkt $(n,k)$
$\binom nk2^k$ Möglichkeiten. Addition über die getrennten, vollständigen
Endpunktklassen liefert
\[
 \sum_{k=0}^n\binom nk2^k=3^n.
\]
Bei $n=0$ besitzt der leere Weg genau eine leere Färbung:
links $\binom00\,2^0=1$, rechts $3^0=1$.
Auch der binomische Lehrsatz mit $a=2$, $b=1$ liefert unmittelbar
$(2+1)^n=\sum_{k=0}^n\binom nk2^k1^{n-k}=\sum_{k=0}^n\binom nk2^k$.

\Needspace{7\baselineskip}
\subsection*{Zu Aufgabe 6: Der letzte Rechtsschritt}
\textbf{(a)} Vor dem letzten $R$ liegen genau zwei $R$ und $i$ $L$.
Der Anfangsweg endet somit bei $(i+2,2)$; es gibt $\binom{i+2}2$
Möglichkeiten. Anschließend kommen das letzte $R$ und genau $n-i$
Schritte $L$. Dieser Rest ist eindeutig festgelegt.

Umgekehrt ergänzen wir jeden Weg nach $(i+2,2)$ durch ein $R$ und
$n-i$ Zeichen $L$. So entsteht genau ein Weg nach $(n+3,3)$, dessen
letztem $R$ genau $i$ Linksschritte vorangehen. Abschneiden und
Ergänzen sind umkehrbar. Jeder vollständige Weg gehört zu genau
einer Klasse $i\in\{0,\ldots,n\}$. Deshalb
\[
 \sum_{i=0}^n\binom{i+2}2=\binom{n+3}3.
\]

\textbf{(b) Induktionsanfang.} Bei $n=0$ steht $\binom22=\binom33$;
beide Werte sind $1$. Im Wegemodell ist dies der einzige Weg $RRR$.

\textbf{Induktionsschritt.} Die Formel gelte für ein $n\in\N$. Dann
\[
 \begin{aligned}
 \sum_{i=0}^{n+1}\binom{i+2}2
 &=\left(\sum_{i=0}^n\binom{i+2}2\right)+\binom{n+3}2\\
 &=\binom{n+3}3+\binom{n+3}2\\
 &=\binom{n+4}3.
 \end{aligned}
\]
Die zweite Gleichheit nutzt die Induktionsannahme, die dritte Pascal.
Weil $n+4=(n+1)+3$, ist dies die Behauptung für $n+1$; die Formel
gilt daher für alle $n\in\N$.

\textbf{(c)} Die zehn Zahlen sind $\binom22,\binom32,\ldots,\binom{11}2$.
Somit ist $n=9$, und die Summe beträgt $\binom{12}3=220$, wie in Aufgabe~3.

\Needspace{7\baselineskip}
\subsection*{Zu Aufgabe 7: Trennen an einer Zwischenzeile}
\textbf{(a)} Eine Station in Zeile $2$ hat die Form $(2,k)$.
Zum Ziel bleiben drei Schritte, davon $3-k$ nach rechts:
\begin{center}
\renewcommand{\arraystretch}{1.25}
\begin{tabular}{@{}ccccc@{}}
\toprule
Station & $k$ & Hinwege & Weiterwege & Ganze Wege\\
\midrule
A & $0$ & $\binom20=1$ & $\binom33=1$ & $1\cdot1=1$\\
B & $1$ & $\binom21=2$ & $\binom32=3$ & $2\cdot3=6$\\
C & $2$ & $\binom22=1$ & $\binom31=3$ & $1\cdot3=3$\\
\bottomrule
\end{tabular}
\end{center}
Insgesamt erhalten wir $1+6+3=10=\binom53$ Wege.

\textbf{(b)} An jeder Station passt jeder Hinweg zu jedem Weiterweg.
Durch Zusammenfügen entsteht ein ganzer Weg; durch Trennen nach zwei
Schritten gewinnen wir die beiden Teile eindeutig zurück. Deshalb
multiplizieren wir die Anzahlen. Jeder Weg trifft Zeile $2$ an genau
einer Station. Die drei Stationsklassen sind damit getrennt und
vollständig; deshalb addieren wir ihre Anzahlen.

Hier ist $m=2$ die Länge des ersten und $n=3$ die Länge des zweiten
Teils; insgesamt sind $r=3$ Rechtsschritte nötig. In Vandermonde lautet
die Summe $\sum_{k=0}^3\binom2k\binom3{3-k}$.
Der zusätzliche Summand $k=3$ ist $\binom23\binom30=0$.

Der Faktor $\binom3{3-k}$ zählt \emph{lokale Weiterwege} ab der
Station: drei verbleibende Schritte mit $3-k$ Rechtsschritten.
Er ist nicht die Anzahl aller Wege von der ursprünglichen Spitze
zu einem späteren Punkt.

\textbf{(c)} Sechs Rechtsschritte sind in fünf Schritten unmöglich.
Auch jeder Summand ist $0$: Für $k>2$ verschwindet $\binom2k$;
für $k\le2$ ist $6-k\ge4>3$, also $\binom3{6-k}=0$.
Damit sind sämtliche $k=0,\ldots,6$ erfasst, und die Summe ist $0$.

\Needspace{7\baselineskip}
\subsection*{Zu Aufgabe 8: Die Mitte und die Spiegelung}
\textbf{(a)} Zur Station $(4,k)$ führen $\binom4k$ Hinwege. Danach
bleiben vier Schritte mit $4-k$ Rechtsschritten, also
$\binom4{4-k}=\binom4k$ Weiterwege nach der Symmetrie.
Jeder Hinweg passt zu jedem Weiterweg; der Schnitt gewinnt beide
Teile eindeutig zurück. Der Beitrag der Station ist somit $\binom4k^2$.
Jeder Weg trifft genau eine Station, also
\[
 \sum_{k=0}^4\binom4k^2=1+16+36+16+1=70=\binom84.
\]
Die rechte Seite zählt alle Wege zum gemeinsamen Ziel $(8,4)$.

\textbf{(b)} Zeile $3$, rechts mit einer Null fortgesetzt, ist
$1,3,3,1,0$. Daher
\[
 1\cdot3+3\cdot3+3\cdot1+1\cdot0=3+9+3+0=15=\binom64.
\]

\textbf{(c)} Sei zunächst $n\ge1$. Ein Weg nach $(2n,n+1)$ hat $2n$
Schritte, darunter $n+1$ Rechtsschritte. Enthält seine erste Hälfte
$k$ Linksschritte, so enthält sie $n-k$ Rechtsschritte. Die Anzahl
der möglichen ersten Hälften ist nach der Spiegelung
\[
 \binom n{n-k}=\binom nk.
\]
In der zweiten Hälfte fehlen noch $(n+1)-(n-k)=k+1$ Rechtsschritte.
Dafür gibt es $\binom n{k+1}$ Möglichkeiten. Jeder erste Teil lässt
sich mit jedem passenden zweiten Teil verbinden; das Trennen nach
$n$ Schritten kehrt diese Verbindung um. Der Beitrag dieser Klasse
ist daher $\binom nk\binom n{k+1}$.

Jeder vollständige Weg hat eine eindeutige Zahl $k$ von Linksschritten
in seiner ersten Hälfte. Die Klassen sind getrennt und erfassen
sämtliche Wege zum Ziel. Für $k=n$ wäre die zweite Hälfte unmöglich,
weil sie $n+1$ Rechtsschritte enthalten müsste; ihr Beitrag ist $0$.
Addition über $k=0,\ldots,n$ ergibt daher
\[
 \sum_{k=0}^n\binom nk\binom n{k+1}=\binom{2n}{n+1}.
\]
Für $n=0$ ist links $\binom00\binom01=0$ und rechts $\binom01=0$.
Ein Schritt nach rechts bei insgesamt null Schritten ist unmöglich.
